\section{A finite candidate search with explicit abstention}
\label{s4:organization-algorithm}
An initial automatic organization operation can be specified without
claiming to infer arbitrary economic meaning. Let a candidate observation
plan be $c=(k,e,a,r,g,u)$: record key $k$, event and availability clocks
$e,a$, revision relation $r$, aggregation $g$, and entity map/unit $u$.
The input includes a finite grammar $\mathcal C$, raw versioned records,
decision identifiers and times, declared units, and constraints with their
provenance. The grammar is itself part of the scientific input. A method
cannot recover an interpretation that the grammar excludes.

The implemented reference grammar has two timestamp assignments, two
record-key hypotheses (source/record or source/entity/record), latest-version
versus append semantics, stock versus flow aggregation, and account versus
a supplied candidate household map. This gives 32 initial candidates.
The engine tests key uniqueness, finite values, currency compatibility,
declared column roles, and any supplied chronology, revision-publication
or accounting constraints. For example, $e\le a$ can distinguish two
timestamp assignments only if the analyst has declared that this feed
records realized events before their publication. Scheduled future events
do not obey that restriction. An external closing-minus-opening ledger
identity can distinguish flow summation from a last stock only if that
identity and its independent measurement are actually supplied.

For each surviving candidate and each requested decision time $t$, first
keep versions with $a<t$. Within its record identity, select the largest
eligible version; then restrict its economic date to the requested
half-open interval $[b,t)$. A correction moving an event outside the
interval must not resurrect the older version. Distinct decision identifiers
and distinct events remain distinct even when their amounts match.
A flow is the sum of distinct eligible increments. A stock is its latest
eligible economic observation; component accounts at that instant can be
summed only under the supplied additive household interpretation. Missing
support returns a missing value, not a fabricated zero. The output retains
the source/entity/record/version identifiers of every selected term.

Let $V_c$ be this materialized view including units and lineage. Hard
constraints define $\mathcal C_0\subseteq\mathcal C$. A transparent
inspection score is
\[
 s(c)=\sum_j\ind\{\text{candidate field }j\text{ agrees with supplied hint }j\}.
\]
Hints order inspection; they cannot eliminate a surviving material
alternative. Partition the survivors by equality of requested values,
lineage, aggregation meaning and economic unit. Return incompatible if
none survive, accept the current requested view if one class remains,
and abstain from a unique view otherwise. Even accepted equality for the
current decisions does not identify a rule for unseen data. Changing the
request can separate previously equivalent plans.

This construction has a simple conditional correctness statement. If the
true plan belongs to the grammar and every hard constraint is valid for
it, that plan survives. If all surviving plans agree on the required view,
the accepted view equals its oracle counterpart. No probabilistic calibration
is needed for this implication. Incorrect metadata or an incomplete grammar
invalidates its premise. A large hint score cannot estimate the probability
that those premises are true. Under ambiguity, fit the scientific functional
separately in each view and retain the resulting sensitivity set; do not
average the tables or treat hint scores as linkage probabilities.

\subsection{An evaluation that reaches the effect estimate}
The synthetic experiment fixes the oracle before constructing candidates.
For 240 customers draw independent standard-normal $X,W$, assign an offer
with probability $\operatorname{expit}(.7X)$, and generate
$Y=20A+30X+\epsilon$, $\epsilon\sim N(0,25)$. The predecision record
contains $X$; a later revision contains $X+2A$ and is unavailable at the
decision. OLS on $(1,A,X)$ estimates the constant effect of 20 in this
specified law. Using the eventual revision instead gives the population
equation $Y=-40A+30(X+2A)+\epsilon$: the nominal cleaning choice changes
the adjustment variable into a treatment descendant.

Across 300 repetitions, labeled revisions and the valid chronology case
accept the oracle view in every run, with effect RMSE about $0.68$--$0.69$.
Retrospective revision adjustment has RMSE about $60$. In the unlabeled
aggregation case an earlier independent record $W$ leaves last-stock $X$
and flow-sum $X+W$ as materially different possibilities. In the linkage
case pairing two accounts changes the available confounder to their sum.
Both cases abstain in every run and retain the oracle. Forcing the top
hint-ranked candidate gives RMSE $10.19$ and $10.59$, while the mean
spread between alternative estimates is $9.77$ and $10.13$. Those spreads
are observation sensitivity, not confidence intervals.

An accounting case has increments $W$ and $X-W$, with an independently
declared reconciled total $X$. The constraint selects the flow plan in every
run, with RMSE $0.70$. The additional total provides the distinguishing
information; the algorithm did not discover it from column names.
Zero wrong acceptances among these fixtures is conditional evidence for
this grammar and generator. It does not establish population error rates
for previously unseen bank feeds. With zero errors in 300 trials, the
two-sided 95\% Wilson upper bound is about $1.27\%$ for the particular
repeated experiment, not zero.

The bounded search costs $O(|\mathcal C|RD)$ in a direct implementation
with $R$ records and $D$ decisions. Indexing by entity and availability
time, caching common candidate projections and expressing eligible-version
selection in a columnar engine reduce repeated scans without changing the
semantics. Probabilistic record linkage is a separate extension requiring
a calibrated joint linkage/outcome model or an explicitly conservative
linkage set. The practical stopping point is a reproducible observation
plan or a short list of decision-relevant alternatives that an analyst can
resolve with additional metadata or measurement.
